\clearpage
\VerdierPhysicalPageAnchor{0234}%
\begingroup\small\setlength{\parskip}{0pt}\setlength{\emergencystretch}{1.5em}
\setlength{\abovedisplayskip}{4pt plus 2pt minus 2pt}%
\setlength{\belowdisplayskip}{4pt plus 2pt minus 2pt}%

\VerdierSemanticAnchor{ch4d-4-1-9-cont}
\noindent Denote by
$\bigl({}^{1}\!Y(p,q),{}^{1}\!\delta(p,q,r)\bigr)$ (respectively,
$\bigl({}^{2}\!Y(p,q),{}^{2}\!\delta(p,q,r)\bigr)$) the spectral object
with values in $K(\mathcal A)$ defined by the complex $X$ filtered by its
first (respectively, second) degree (4.1.7). It follows immediately from
the definitions that one has the relations

\VerdierSemanticAnchor{ch4d-4-1-9-1}
\begin{equation}
\left\{
\begin{aligned}
{}^{1}\!Y(q-1,q)&=X^{-q,\bullet}[q],\\
{}^{2}\!Y(q-1,q)&=X^{\bullet,-q}[q],\\
{}^{1}\!\delta(q-1,q,q+1)&=-\bigl(d_1^{-q-1,\bullet}\bigr)[q+1],\\
{}^{2}\!\delta(q-1,q,q+1)&=-\bigl(d_2^{\bullet,-q-1}\bigr)[q+1].
\end{aligned}
\right.
\tag{4.1.9.1}
\end{equation}
These formulas will be used in (4.6).

\VerdierSemanticAnchor{ch4d-4-2}
\medskip\noindent\textbf{4.2. Spectral object with values in
$D(\mathcal A)$ associated with a filtered object of
$\operatorname{comp}(\mathcal A)$.}

\VerdierSemanticAnchor{ch4d-4-2-scope}
\medskip\noindent\emph{In this subsection and the following subsections of
this section, $\mathcal A$ denotes an abelian category.}

\VerdierSemanticAnchor{ch4d-4-2-1}
\medskip\noindent\textbf{4.2.1.}\quad
Let $X$ be an object of $\operatorname{comp}(\mathcal A)$ and let
$X(q)$, $q\in\mathbb Z$, be an increasing filtration of type
$\widetilde{\mathbb Z}$ of the complex $X$:

\VerdierSemanticAnchor{ch4d-4-2-1-1}
\begin{equation}
0=X(-\infty)\subset\cdots\subset X(q)\subset X(q+1)
\subset\cdots\subset X(+\infty)=X.
\tag{4.2.1.1}
\end{equation}

\VerdierSemanticAnchor{ch4d-4-2-1-associated-object}
The construction recalled in (4.1.1) yields a spectral object
$\bigl(\mathcal X(p,q),\delta(p,q,r)\bigr)$ with values in
$K(\mathcal A)$ and hence, on applying the canonical functor
$K(\mathcal A)\longrightarrow D(\mathcal A)$, a spectral object with
values in $D(\mathcal A)$, still denoted, by an abuse of notation, by
$\bigl(\mathcal X(p,q),\delta(p,q,r)\bigr)$. We shall describe a spectral
object with values in $D(\mathcal A)$ isomorphic to this last spectral
object.

\VerdierSemanticAnchor{ch4d-4-2-2}
\medskip\noindent\textbf{4.2.2.}\quad
Thus let
$0=X(-\infty)\subset\cdots\subset X(q)\subset\cdots\subset
X(+\infty)=X$ be a filtered object of $\operatorname{comp}(\mathcal A)$.
Set

\VerdierSemanticAnchor{ch4d-4-2-2-1}
\begin{equation}
X(p,q)=X(q)/X(p),
\tag{4.2.2.1}
\end{equation}
for every element $(p,q)$ of
$\mathcal{F}\ell_1(\widetilde{\mathbb Z})$. (The quotient is of course
taken here in the sense of the abelian category
$\operatorname{comp}(\mathcal A)$.) Denote by

\VerdierSemanticAnchor{ch4d-4-2-2-2}
\begin{equation}
\varphi\bigl((p,q),(p',q')\bigr)\colon
X(p,q)\longrightarrow X(p',q'),
\tag{4.2.2.2}
\end{equation}
\endgroup
